\section{引言}

低照度图像增强在夜间拍摄、视频监控、移动终端成像和机器视觉等场景中具有重要价值。受曝光不足、噪声放大和局部对比度偏低等因素影响，原始低照度图像常存在暗区细节难以辨认、纹理信息不明显和视觉观感较差等问题。因此，设计一类既能提升局部对比度、又能抑制暗区噪声和块效应的增强方法，是数字图像处理中的一个典型问题。

现有低照度增强方法大致可以分为三类：以 HE、AHE、CLAHE 为代表的传统对比度增强方法\cite{pizer1987ahe,zuiderveld1994clahe}，以 Retinex、照明估计和图像融合为代表的照明建模方法\cite{jobson1997retinex,wang2013loe,fu2016fusion,guo2017lime}，以及兼顾亮度恢复与噪声抑制的联合增强方法\cite{li2015lowlight,ren2018joint}。其中，直方图均衡及其改进方法由于无需训练、实现简单、可解释性强，仍然是数字图像处理课程背景下最适合深入分析的一条经典路线。CLAHE（Contrast Limited Adaptive Histogram Equalization）通过局部直方图均衡和对比度限制在很多场景中取得了比全局均衡更好的视觉效果\cite{pizer1987ahe,zuiderveld1994clahe}。然而，固定参数 CLAHE 仍存在两个显著不足：其一，平坦暗区往往会被过增强，从而放大噪声；其二，局部参数突变和硬拼接容易引起块状伪影。

围绕上述问题，本文并不试图提出一种全新的低照度增强范式，而是以“改进 CLAHE 的 clip limit 生成方式”为核心切入点，提出一种噪声抑制的局部方差自适应 CLAHE 方法。与直接利用原始局部方差的方法不同，本文使用平滑亮度图上的结构方差并结合残差域 MAD 噪声估计构造结构强度图，在此基础上自适应分配 tile 级 clip limit；此外，为降低块边界突变，本文进一步引入 tile 间双线性插值融合。本文工作的主要贡献概括如下：

\begin{enumerate}
  \item 提出基于残差域噪声估计的结构方差构造方式，避免噪声被直接误判为结构纹理；
  \item 设计分位数归一化的自适应 clip limit 映射和 tile 间双线性融合机制，用于减弱硬拼接带来的块效应；
  \item 在合成低照度、真实低照度、消融实验、融合策略对照和更强退化条件下进行了系统验证，并对方法的局限性进行讨论。
\end{enumerate}
